%======================================================================
%  Workload concepts and analytical model
%======================================================================

\section{A performance model for LLM inference}
\label{sec:inference-model}

\subsection{Prefill and autoregressive decode}

Inference for a decoder-only transformer has two phases
~\cite{na2024llmcpu}.  During
\emph{prefill}, the model processes the prompt and constructs the
initial key--value (KV) cache.  Tokens from the prompt can be processed
in parallel, so the linear layers often form comparatively large matrix
multiplications.  During \emph{decode}, the model generates one new
token per sequence at each iteration.  Each iteration requires the
relevant weights and KV state, then appends the key and value for the new
token.  Dense full-context attention reads the prior cached KV; sliding
window or sparse attention can reduce that history.

\begin{figure}[tbp]
\centering
\begin{tikzpicture}[
  font=\small,
  io/.style={draw=black!50, fill=white, rounded corners=2pt,
             minimum width=2.2cm, minimum height=0.9cm, align=center},
  phase/.style={draw=black!65, fill=neutralcol, rounded corners=3pt,
                minimum width=4.0cm, minimum height=1.55cm, align=center},
  kv/.style={draw=hbmcol!85!black, fill=hbmcol!16, rounded corners=2pt,
             minimum width=4.0cm, minimum height=0.8cm, align=center},
  flow/.style={-{Stealth}, thick},
  note/.style={font=\footnotesize, align=center}
]
  \node[io] (prompt) at (-6.0,0) {prompt};
  \node[phase] (prefill) at (-2.5,0)
    {\textbf{Prefill}\\process prompt in parallel\\build initial KV};
  \node[phase] (decode) at (2.6,0)
    {\textbf{Decode loop}\\access weights and relevant KV\\
     generate one token/sequence};
  \node[io] (output) at (6.3,0) {token stream};
  \node[kv] (cache) at (2.6,-2.1)
    {KV cache: read required history, append new state};

  \draw[flow] (prompt) -- (prefill);
  \draw[flow] (prefill) -- node[above, note]{first token} (decode);
  \draw[flow] (decode) -- (output);
  \draw[flow, {Stealth}-{Stealth}] (decode) -- (cache);
  \draw[flow] (decode.45) .. controls +(0.8,0.8) and +(-0.8,0.8)
       .. (decode.135) node[midway, above, note]{next iteration};

  \node[note, below=0.2cm of prefill]
    {principal latency metric: TTFT};
  \node[note, below=0.65cm of cache]
    {principal latency metric: inter-token latency};
\end{tikzpicture}
\caption{Dataflow of prompt prefill and autoregressive decode.  The
frequent description of prefill as compute-bound and decode as
bandwidth-bound is a useful first-order model, not a universal property:
batch, context, precision and kernel implementation can change the
regime.  Author's synthesis.}
\label{fig:prefill-decode}
\end{figure}

This phase distinction is essential when evaluating HBM.  A single
end-to-end throughput number can hide a faster decode phase behind a
long prompt, or hide a prefill regression behind a long generation.
TTFT and decode latency must therefore be reported separately.  Prefill
usually offers matrix--matrix parallelism, whereas low-batch decode is
closer to matrix--vector execution and repeatedly streams weights.
Larger batches increase weight reuse and can move decode toward
matrix-engine throughput; long contexts simultaneously increase KV
traffic.

\FloatBarrier

\subsection{Working-set capacity}
\label{sec:working-set}

For a model with \(P\) parameters and a stored weight precision of
\(q_w\) bits, the first-order weight footprint is

\[
  M_W = P\frac{q_w}{8} + M_{\mathrm{format}},
\]

where \(M_{\mathrm{format}}\) includes scales, zero-points, alignment,
and any prepacked or duplicated representation created by the runtime.
The model file size is not necessarily the resident weight footprint.

For \(B\) concurrent sequences, cached length \(S\), \(L\) transformer
layers, \(H_{\mathrm{KV}}\) KV heads, head dimension \(d_h\), and
\(q_{\mathrm{KV}}\) bits per cached element, the KV footprint is

\[
  M_{\mathrm{KV}} =
  2 B S L H_{\mathrm{KV}} d_h \frac{q_{\mathrm{KV}}}{8}.
\]

The factor two accounts for keys and values.  Grouped-query and
multi-query attention reduce the KV footprint because
\(H_{\mathrm{KV}}\) is smaller than the number of query heads.  The
complete resident set is

\[
  M_{\mathrm{live}} =
  M_W + M_{\mathrm{KV}} + M_{\mathrm{activations}}
  + M_{\mathrm{workspace}} + M_{\mathrm{runtime}}.
\]

This arithmetic identifies where capacity boundaries may occur, but the
boundary must use the OS-exposed NUMA capacity on the allocated node and
a measured safety margin.  The advertised 64\,GB per socket is not a
usable-allocation guarantee: system use, allocator fragmentation and
workspace consume part of it.  A node may expose roughly 128\,GB across
two sockets, but that capacity is not uniformly local; using it requires
a placement strategy that distributes data and execution deliberately.

Quantization changes both sides of the experiment.  Fewer bits reduce
weight capacity and memory traffic, potentially reducing the benefit of
HBM for a model that already fits.  At the same time, quantization can
move a larger model below the HBM capacity boundary and remove DDR or
remote traffic altogether.  Packing, dequantization and kernel coverage
prevent performance from scaling exactly with bit width.  Weight and KV
precision must therefore be reported separately.

\subsection{Arithmetic intensity and the Roofline bound}
\label{sec:roofline}

The Roofline model relates arithmetic intensity \(I=F/D\), measured in
FLOP/byte, to compute throughput and memory bandwidth
~\cite{williams2009roofline}:

\[
  P_{\mathrm{attainable}}
  \leq
  \min\!\left(P_{\mathrm{compute}},
              I B_{\mathrm{effective}}\right).
\]

For one inference step, a corresponding lower-bound model is

\[
  T_{\mathrm{step}}
  \gtrsim
  \max\!\left(
      \frac{F_{\mathrm{step}}}{P_{\mathrm{effective}}},
      \frac{D_{\mathrm{step}}}{B_{\mathrm{effective}}}
  \right)
  + T_{\mathrm{communication}} + T_{\mathrm{runtime}}.
\]

The effective quantities must be measured for the selected kernel,
placement and core count.  Processor peak FLOP/s and theoretical pin
bandwidth are ceilings, not values to substitute uncritically.

For a dense linear layer during decode, if a stored weight is used once
for each of \(B\) sequences in a batch, a useful first-order
approximation is

\[
  I_{\mathrm{weights}} \approx \frac{2B}{b_w}
  \quad \text{FLOP/byte},
\]

where \(b_w\) is the stored bytes per weight and a multiply--add is
counted as two FLOPs.  The approximation assumes that each fetched dense
weight is reused across the \(B\) sequences.  At batch one, BF16 weight
streaming is therefore close to one FLOP/byte before KV,
activation and runtime traffic are included.  Batching reuses weights
and raises arithmetic intensity, which explains why HBM is expected to
matter most for low-batch decode and why AMX can become more important
as the matrix dimensions grow.

This model also clarifies why bandwidth alone is insufficient:

\begin{itemize}
  \item \textbf{Bandwidth} sets the rate for many concurrent,
        sufficiently regular memory requests.
  \item \textbf{Latency} sets the cost of a dependent or lightly loaded
        access path.  HBM can have higher unloaded latency than DDR5
        while sustaining much greater bandwidth under load
        ~\cite{ibeid2025aurora}.
  \item \textbf{Capacity} determines whether the critical working set
        remains local.  Crossing a 64\,GB socket boundary can introduce
        DDR5 or UPI traffic and cause a discontinuity that a simple
        single-tier Roofline does not predict.
\end{itemize}

For concurrent HBM and DDR5 traffic, a no-contention service-time
estimate can range from perfect overlap to serialization of the
per-tier transfers.  It is not a general bound: shared uncore resources,
cache conflicts, UPI traffic and memory-controller interference can make
the measured time worse than the serialized estimate.  Simultaneous
HBM--DDR traffic must therefore be characterized before a tiered
placement policy is modelled.

\subsection{Placement is part of the algorithm}

In Flat mode the operating system exposes memory tiers through NUMA,
but it does not know which model objects are bandwidth-critical.
Several mechanisms can therefore produce the wrong placement:
allocator first touch, worker migration, memory-mapped model pages,
automatic NUMA balancing, and a runtime that creates prepacked copies
after the initial allocation.  A command-line memory policy is evidence
of intent, not evidence of residency.

The relevant experimental object is consequently not ``the model in
HBM'' but a verified mapping between:

\begin{itemize}
  \item worker threads and physical cores;
  \item weight, KV, activation and workspace pages and their NUMA nodes;
  \item local or remote memory traffic observed during each inference
        phase; and
  \item the time interval used for performance and energy measurement.
\end{itemize}

This mapping is also the basis for testing topology-aware execution.
Independent replicas can use one local socket each; a single model can
be sharded across sockets or SNC domains; and selected objects can be
placed in different memory tiers.  These designs answer different
questions and should not be grouped under a generic ``two-socket run''.

\subsection{Metrics and boundaries}

For \(O\) generated tokens in a synchronous request,

\[
  T_{\mathrm{E2E}}
  =
  T_{\mathrm{TTFT}} +
  \sum_{j=2}^{O} T_{\mathrm{ITL},j}.
\]

At minimum, results should report TTFT, median and tail ITL, end-to-end
latency, output tokens/s and requests/s.  The report must state whether
tokenization, model loading, queueing and warm-up are included.  Static
batch throughput and online serving goodput are different constructs:
the latter additionally depends on arrivals, scheduling and declared
latency objectives.

Hardware counters are explanatory measurements rather than performance
metrics.  HBM and DDR traffic, UPI traffic, memory-controller
utilization, cache misses, achieved frequency and AMX-related counters
can support a mechanism only when their collection overhead,
multiplexing and scope are documented.

\subsection{Predictions to test}

The model yields three bounded predictions:

\begin{enumerate}
  \item local HBM should benefit low-batch decode more than compute-dense
        prefill, but the speedup need not equal the STREAM bandwidth
        ratio;
  \item performance should change when the verified live set crosses a
        local HBM boundary, with the size and direction of the change
        depending on placement and memory mode; and
  \item batching and lower-bit weights should reduce weight-bandwidth
        pressure, while long contexts can reintroduce pressure through
        the KV cache.
\end{enumerate}

These statements are hypotheses.  The literature reviewed next
motivates them but does not establish them for CRESCO8.
